\documentclass[14pt,a4paper]{extarticle} % Permite 8pt, 9pt, 10pt, 11pt, 12pt, 14pt, 17pt, 20pt
\usepackage[utf8]{inputenc}
\usepackage[spanish]{babel}
\usepackage[margin=1.5cm]{geometry}
\usepackage{pdflscape}
\usepackage{tabularx}
\usepackage{booktabs}
\usepackage{ragged2e}
\usepackage{enumitem}
\usepackage[table]{xcolor}
\usepackage{setspace} % Para controlar el interlineado de las líneas de texto

% --- CONFIGURACIÓN DE ESPACIADO ---
\renewcommand{\arraystretch}{1.8} % Espacio entre el texto y las líneas de la celda
\setlength{\tabcolsep}{10pt}      % Espacio entre columnas

% Definición de columna Y con interlineado de 1.1 y fuente más pequeña
\newcolumntype{Y}{>{\RaggedRight\setstretch{1.1}\footnotesize\arraybackslash}X}

\begin{document}

\begin{landscape}
\thispagestyle{empty}
\begin{center}
    \textbf{\Large MATRIZ DE CONSISTENCIA} \\
    \vspace{0.2cm}
    \small Proyecto: Análisis comparativo del rendimiento computacional de algoritmos (IA vs Estándar)
\end{center}

\vspace{0.3cm}

\noindent
\begin{tabularx}{\linewidth}{|Y|Y|Y|Y|Y|}
\hline
\rowcolor[gray]{0.9} 
\textbf{Problemas} & \textbf{Objetivos} & \textbf{Hipótesis} & \textbf{Variables} & \textbf{Metodología} \\ \hline

\textbf{Problema General:} \newline
¿Qué diferencias existen en el rendimiento computacional, medido en tiempo de ejecución y consumo de memoria, entre los algoritmos generados por IA generativa y las implementaciones estándar? \newline

\textbf{Problemas Específicos:}
\begin{itemize}[leftmargin=*, nosep, itemsep=6pt] % Mayor separación entre puntos
    \item \textbf{PE1:} ¿Diferencias en tiempo de ejecución bajo diferentes tamaños de entrada?
    \item \textbf{PE2:} ¿Variación del consumo de memoria máxima (Peak RSS) ante distintos volúmenes de datos?
    \item \textbf{PE3:} ¿Influencia del escalado de datos en la brecha de rendimiento?
\end{itemize} 
& 
\textbf{Objetivo General:} \newline
Analizar comparativamente el rendimiento computacional de algoritmos de ordenación y grafos (IA vs Estándar) evaluando tiempo y memoria. \newline

\textbf{Objetivos Específicos:}
\begin{itemize}[leftmargin=*, nosep, itemsep=6pt]
    \item \textbf{OE1:} Evaluar el tiempo de ejecución bajo diferentes tamaños de entrada.
    \item \textbf{OE2:} Analizar el consumo de memoria máxima ante distintos volúmenes de datos.
    \item \textbf{OE3:} Determinar el impacto del escalado de datos en la degradación del rendimiento.
\end{itemize} 
& 
\textbf{Hipótesis General:} \newline
Existe una diferencia significativa; los algoritmos por IA presentan menor eficiencia en tiempo y memoria frente a las implementaciones estándar. \newline

\textbf{Hipótesis Específicas:}
\begin{itemize}[leftmargin=*, nosep, itemsep=6pt]
    \item \textbf{HIP1:} El tiempo de ejecución en IA es significativamente mayor.
    \item \textbf{HIP2:} El consumo de memoria es superior en IA por ineficiencias en gestión de estructuras.
    \item \textbf{HIP3:} La brecha de ineficiencia se acentúa de forma no lineal al aproximarse a $10^6$ elementos.
\end{itemize} 
& 
\textbf{Variables Independientes:}
\begin{itemize}[leftmargin=*, nosep, itemsep=2pt]
    \item \textbf{V.I.1:} Origen del Código.
    \item \textbf{V.I.2:} Tamaño de Entrada ($n$).
    \item \textbf{V.I.3:} Tipo de Algoritmo.
\end{itemize}

\textbf{Variables Dependientes:}
\begin{itemize}[leftmargin=*, nosep, itemsep=2pt]
    \item \textbf{V.D.1:} Tiempo de Ejecución.
    \item \textbf{V.D.2:} Consumo de Memoria.
\end{itemize} 
& 
\textbf{Método:} Experimental / Cuantitativo. \newline
\textbf{Tipo:} Aplicada / Tecnológica. \newline
\textbf{Diseño:} Experimental transversal. \newline

\textbf{Muestra:} \newline
QuickSort, MergeSort y Dijkstra (IA vs Estándar). \newline

\textbf{Técnicas:} Benchmarking. \newline
\textbf{Instrumentos:} Hyperfine y GNU time. \\ \hline
\end{tabularx}

\end{landscape}

\end{document}